% A proof contribution for the parent article; not a standalone deliverable.
\section{The exact rectangular cancellation criterion}
\label{sec:column-peeling}

The maintained core solver first peels rows and, for rectangular matrices,
falls back to a saturated positive-pivot heuristic.  There is a dual source
criterion that decides full row completion directly.  This is a consequence
of report~56's raw-epoch theorem\cite{raw-epoch-56}, applied to the transposed incidence matrix
with an empty survivor alphabet.  We give a direct proof because it reveals
a short obstruction certificate and avoids searching over surviving columns.
No priority claim is made for the underlying matching or closure criterion.

Let $A=(a_{ic})$ be a nonnegative integral $m\times n$ matrix.  A
\emph{positive unit pivot} at $(r,c)$ is allowed when $a_{rc}=1$; it deletes
row $r$ and column $c$ and changes every surviving entry to
\[
 a'_{ih}=a_{ih}+a_{ic}a_{rh}.
\]
The question is whether some sequence deletes every row.  All occurrences
are unsigned geometric incidences; signed exponent sums and freely reduced
words define a different problem.

A \emph{column peel} deletes row $r$ using a column $c$ whose total over
the live rows is exactly one.  Then $a_{rc}=1$ and $a_{ic}=0$ for $i\ne r$,
so the positive update has no fill.  A set $U$ of source rows is a
\emph{stopping set} when
\[
              \sum_{i\in U}a_{ic}\ne1\qquad\text{for every column }c.
\]
The empty set is allowed, and a nonempty zero row is itself an obstruction.

\begin{lemma}[Reverse unit-pivot lemma, recalled with proof]
\label{lem:reverse-positive-pivot-self-contained}
Let $A$ be square and let $B$ result from a positive unit pivot.  If
$\operatorname{per}(B)=1$, then $\operatorname{per}(A)=1$.
\end{lemma}
\begin{proof}
Regard each entry as the number of separately labelled edges between its row
and column.  An edge counted by $b_{ih}$ has a tag indicating either an
original edge counted by $a_{ih}$, or an ordered pair consisting of an
$a_{ic}$ edge and an $a_{rh}$ edge.  Call the latter a composite edge.

The unique tagged perfect matching of $B$ has at most one composite edge.
Otherwise, take two such edges $(i,h)$ and $(j,k)$.  Exchanging their second
factors gives valid composite edges $(i,k)$ and $(j,h)$, and hence a distinct
perfect matching.  If there is no composite edge, append the unique pivot
edge $(r,c)$ to obtain a matching of $A$.  If there is one, replace it by
its two labelled source edges.  This again produces a matching of $A$.

There is also an injection from perfect matchings of $A$ to tagged perfect
matchings of $B$.  A matching using $(r,c)$ maps to its restriction, with all
edges tagged as original.  A matching not using $(r,c)$ has unique edges
$(i,c)$ and $(r,h)$; replace them by their composite edge $(i,h)$ and tag the
remaining edges as original.  The image, including its tags, recovers the
source matching in either case.  Therefore $A$ has at most one matching,
and the preceding construction shows that it has at least one.
\end{proof}

\begin{theorem}[Full-row positive cancellation]
\label{thm:column-peeling}
The following statements are equivalent.
\begin{enumerate}
\item Some positive unit-pivot sequence deletes every row of $A$.
\item $A$ has an $m\times m$ submatrix, using all rows, with permanent one.
\item Column peeling deletes every row, for every choice of eligible column.
\item The source has no nonempty stopping set.
\end{enumerate}
Consequently no enumeration of the $n-m$ surviving columns is required.
\end{theorem}

\begin{proof}
For $m=0$ all statements hold with the empty minor having permanent one.
If a positive-pivot sequence succeeds, record its $m$ chosen source columns.
Restrict to these columns throughout the sequence.  The restricted matrices
obey exactly the same positive-pivot rule and end at the empty matrix.
Lemma~\ref{lem:reverse-positive-pivot-self-contained}, proved in report~56
and recalled above, implies that permanent one after a positive pivot gives
permanent one before it.  Reversing the complete sequence proves the second
statement.  In particular $m\le n$ is necessary.

Suppose a chosen square minor has permanent one.  Its unique edge-labelled
perfect matching uses unit entries.  For any nonempty source-row set $U$,
restrict that matching to $U$ and to its matched columns.  This smaller square
matrix also has permanent one: another matching would extend, using the
unchanged matching outside $U$, to another matching of the original minor.
Relabel its matched entries onto the diagonal.  If every column had another
positive entry, the directed graph of off-diagonal positive entries would
have positive indegree at every vertex and hence a directed cycle.  Exchanging
the matched columns around that cycle would produce a second perfect
matching.  Some column therefore contains just one unit entry in the rows
$U$.  Thus $U$ is not a stopping set, proving the fourth statement.

If peeling stops on a nonempty set $U$, every column has total zero or at
least two on $U$, so $U$ is a stopping set.  Under the fourth statement this
cannot happen, whatever eligible columns were selected.  This proves the
third statement.  Finally every column peel is a legal positive pivot with
zero fill, proving the first statement from the third.
\end{proof}

\begin{theorem}[Canonical obstruction core]
\label{thm:stopping-core}
Column peeling leaves exactly the union $K$ of all source stopping sets.
Thus its residual row set is independent of pivot order.  A nonempty subset
$U$ satisfying the stopping inequalities is already a complete certificate
that full-row positive cancellation is impossible; maximality is unnecessary
for that conclusion.
\end{theorem}
\begin{proof}
The union of stopping sets is a stopping set.  Indeed, if a column has a
positive entry somewhere in the union, that row lies in one constituent
stopping set, where the column total is already at least two.  Its total in
the union is consequently at least two.

No peeling sequence can delete a row of a fixed stopping set $U$: at the
first alleged deletion, every row of $U$ is still present, so a column
meeting the selected row has total at least two.  Therefore the union $K$
survives every peel.  The final residual $R$ is itself a stopping set, hence
$R\subseteq K$.  The two inclusions yield $R=K$.
\end{proof}

\paragraph{The scope of a negative answer.}
The all-one $2\times2$ matrix is a stopping core, but it has four allowed
unit pivots.  Any such pivot leaves the one-entry matrix $(2)$.
Column peeling therefore computes neither a maximum partial cancellation nor
a minimum topological handle decomposition.  An obstructed authenticated
core can still be a handlebody after geometric moves outside this raw model.

\begin{corollary}[Row-subset monotonicity and an early core filter]
\label{cor:row-subset-monotonicity}
If a matrix permits full-row positive cancellation, so does every matrix
obtained by restricting it to a subset of rows.  Equivalently, adding rows
cannot rescue an obstructed source matrix.  This remains true if one also
adds columns whose entries in all original rows are zero.

For the saturated-bundle construction, an original core stopping subset
therefore remains an obstruction in the joint presentation.  With the same
original-core forest fixed, the original core can be checked before
constructing the full ordered bundle boundary interface.  An obstruction
suffices to return an inconclusive result from this particular handlebody
certificate method.
\end{corollary}
\begin{proof}
Take a complete column-peeling schedule and retain only its moves on the
chosen row subset.  At each retained move, the chosen column has the same
single unit occurrence as before: the omitted rows can only reduce its
total, and the selected row itself remains.  Thus the restricted schedule
also completes.

Alternatively, a stopping subset of original rows remains a stopping subset
after adding any other rows.  Its column totals are computed only over the
chosen subset, and any newly added column contributes zero there.  The
saturated compiler leaves original-core rows unchanged in active columns
and gives them zero in private bundle columns.  The same obstruction
therefore applies to its joint matrix.
\end{proof}

This is a statement about the fixed raw geometric presentation.  It is not
the topological assertion that attaching another piece can never turn a
non-handlebody into a handlebody.  A different original-core forest, geometric
normalization, new disk system, or an allowed two/three-handle cancellation
may change the raw problem and lies outside the early filter's conclusion.

\subsection{Sparse discovery and independent verification}

\begin{lemma}[Columns with no unit entries are inert]
\label{lem:inert-columns}
Suppose every initial entry of a column is zero or at least two.  That column
can never contain a unit entry under any positive-pivot sequence.  Removing
any family of such columns leaves the answer to full-row completion unchanged.
The omitted columns must still be counted when computing surviving genus.
\end{lemma}
\begin{proof}
For an omitted column $h$, a surviving entry changes from $a_{ih}$ to
$a_{ih}+a_{ic}a_{rh}$.  If both $a_{ih}$ and $a_{rh}$ are zero or at least
two, their displayed sum is again zero or at least two.  Induction shows that
$h$ can never be a pivot column.  Updates in any other column do not depend
on $h$, so a sequence on the full matrix projects to the same sequence after
omission, and any sequence on the smaller matrix lifts with identical pivot
entries.  Omission is an algebraic search reduction, not deletion of the
corresponding one-handles from the topological object.
\end{proof}

This applies in particular to private columns whose only initial entry is
two, as in the saturated-bundle compiler.  The invariant requires only the
absence of unit entries; divisibility by two is stronger than necessary.

Replace every positive source multiplicity by its cap $\min(2,a_{ic})$.
This leaves all stopping inequalities and column-singleton predicates
unchanged.  For each column maintain the \emph{uncapped sum of these capped
entries} over live rows, together with the xor of their distinct dense row
indices.  A column of total one has exactly one unit occurrence, whose row
is the xor.  On deleting that row, visit its immutable source incidences
once, subtract their capped entries from the column totals, and xor out the
row index.  Enqueue a column when its total first becomes one.

Column totals decrease strictly when visited.  Each column enters the queue
at most once and each nonzero source incidence is removed at most once.  If
$I$ is the number of nonzero source incidences, discovery uses $O(m+n+I)$
expected dictionary operations and working entries, apart from sorting the
reported residual labels.  That final normalization costs $O(m\log(m+1))$
comparisons.  Dense indices and capped-entry column totals have
$O(\log(m+n+1))$ bits.  Arbitrary source-label bit lengths and the time to
read uncapped input integers are charged separately.

The independent verifier does not simulate the queue.  It assigns a deletion
time to every row in the certificate and records the proposed pivot column.
For pivot $(r,c)$ at time $t$ it checks from the immutable source that
$a_{rc}=1$ and that every other row incident to $c$ has a deletion time less
than $t$.  A single scan of source incidences verifies these inequalities for
all moves.  The same scan checks that every residual column has total zero
or at least two.  Thus the certificate proves both a legal deletion sequence
and a fixed point; Theorem~\ref{thm:stopping-core} identifies the latter as
the canonical obstruction.  On a complete result the residual is empty.

\subsection{A correction to the old producer's status description}

The pinned native saturated producer minimizes
\[
  \Bigl(\sum_h a_{rh}-1\Bigr)
  \Bigl(\sum_i a_{ic}-1\Bigr)
\]
over unit pivots, with the entries capped at two.  A zero score means a
unit row leaf or a unit column leaf.  In fact this producer is already
complete for the \emph{full-row} decision, although its docstring describes
it as incomplete.

\begin{proposition}[Completeness of minimum-fill greedy for full completion]
Any positive-pivot producer that always prefers a zero-fill pivot whenever
one exists decides full-row completion correctly, provided it runs until all
rows disappear or no unit pivot remains.
\end{proposition}
\begin{proof}
Assume full completion exists.  Theorem~\ref{thm:column-peeling} supplies a
unit column leaf, so a zero-fill pivot is available.  It remains to show that
either type of zero-fill pivot preserves existence of a permanent-one minor.

For a row leaf, its sole unit column belongs to every matching covering all
rows.  Remove that forced matched pair from a permanent-one minor.  For a
column leaf at row $r$, if the minor already includes that column the same
argument applies.  Otherwise replace the column matched to $r$ by the new
leaf column.  Any matching of the resulting square minor must pair this
column with $r$; another matching on the remaining rows would extend to a
second matching of the old minor.  Thus the replacement minor still has
permanent one.  Removing the leaf pair leaves a permanent-one minor for the
residual matrix.  Since zero-fill pivots do not alter surviving entries,
induction proves completion.

Conversely a completed output is itself a positive-pivot sequence.  If the
source is obstructed, no choices made by any producer can complete it.
\end{proof}

The new solver therefore improves the work bound and the negative certificate;
it does not increase successful inputs merely by replacing that old
minimum-fill routine.  It also avoids the old producer's potentially large
partial fill work once full completion is impossible.

\subsection{Exact work families}

For the $m\times(m+1)$ chain matrix $a_{i,i}=a_{i,i+1}=1$, the new pass
visits $2m$ source incidences and then at most $2m$ removal incidences.
The old producer rebuilds row and column totals at every pivot; its native
\texttt{entry\_scans} counter is $2m(m+1)$ on this family.  The same count
holds for the rectangular fan with $a_{i,i}=1$ and $a_{i,m}\ge2$.
These are source-incidence planner bounds; no crossing-count interpretation
or end-to-end knot-recognition speedup follows from the matrices alone.
